\Askhsh{1727}{4}
\begin{ylh}{1}
    \item 3.1. Είδη και στοιχεία τριγώνων
    \item 3.11. Ανισοτικές σχέσεις πλευρών και γωνιών
    \item 4.2. Τέμνουσα δύο ευθειών - Ευκλείδειο αίτημα
    \item 4.6. Άθροισμα γωνιών τριγώνου
    \item 5.2. Παραλληλόγραμμα
    \item 5.3. Ορθογώνιο
    \item 5.6. Εφαρμογές στα τρίγωνα.
\end{ylh}
Δίνεται τραπέζιο $\mathrm{A}\mathrm{B}\Gamma\Delta$ με $\hat{\mathrm{A}} = \hat{\Delta} = 90^\circ$, $\Delta\Gamma = 2\mathrm{A}\mathrm{B}$ και $\hat{\mathrm{B}} = 3\hat{\Gamma}$. Από το $\mathrm{B}$ φέρνουμε κάθετη στη $\Gamma\Delta$ που τέμνει την $\mathrm{A}\Gamma$ στο σημείο $\mathrm{K}$ και την $\Gamma\Delta$ στο $\mathrm{E}$. Επίσης φέρνουμε την $\mathrm{A}\mathrm{E}$ που τέμνει τη $\mathrm{B}\Delta$ στο σημείο $\Lambda$.
Να αποδείξετε ότι:
\begin{alist}
    \item $\hat{\Gamma} = 45^\circ$ \monades{8}
    \item $\mathrm{B}\Delta = \mathrm{A}\mathrm{E}$ \monades{9}
    \item $\mathrm{K}\Lambda = \frac{1}{4}\Delta\Gamma$ \monades{8}
\end{alist}
